% English translation, made from our own transcription (original.tex), not from any existing
% translation. Every math expression, \origpage marker and the \rmfigure image path are preserved
% unchanged; only the prose is translated. Notes for the reviewer:
%  * The five-ratio table on p. 604 is a single display-math block whose ratio annotations are set
%    inside \text{...} (e.g. "ex nat. Circuli.", "per constr. & princ. Calc. inf. parvor."). Because
%    the math-preservation check (pipeline/texcompare.py) treats the whole display block as one
%    invariant token, those Latin abbreviations are kept VERBATIM inside the block — they cannot be
%    translated without registering as an altered formula. Their sense: "from the nature of the
%    Circle", "by construction & the principles of the calculus of infinitely small quantities",
%    "from the nature of the curve AQ, as gathered from Cor. 2 of Constr. III preceding".
%  * The uncertain reading in note (i), $\uncertain{Ai\ [a] :}$, sits inside a math span and is
%    likewise preserved as-is.
%  * The central differential equation appears in three variant forms in the source (p. 601,
%    note (b), obs. 8), all preserved as printed — a faithful discrepancy, not translated away.
%  * Period abbreviations in the prose are rendered in English (scil. = namely, v. g. = for example,
%    per constr. = by construction, No. praec. = the preceding No.); the author's "&c." is given as
%    "etc.". The lettered notes (a)-(m) keep their inline placement at the foot of each page (R15).
%  * Markup-only fix applied identically to original.tex and this file: at the five sites where a
%    "$...$" span directly abutted a "${}^{(x)}$" footnote marker (f, i, k, l, m), the two spans were
%    merged as "$...\,{}^{(x)}$" — one span with a thin space. This changes neither the mathematics
%    nor the web rendering (the site treats each "$...$" as inline; only "\[...\]" is display), but
%    it removes a literal "$$" that pipeline/texcompare.py would otherwise mis-parse as display math.
\documentclass{article}
\usepackage{readmasters}
\begin{document}

\origpage{601}
\section*{No.~LIX. Jacob Bernoulli's Solution of Leibniz's Problem}
\subsection*{On the Curve of uniform Approach to and Recession from a given point, by means of the rectification of the Elastic Curve}

The felicity of the preceding discovery may be commended by the solution of that most elegant
\emph{Leibnizian} Problem, of finding the \emph{line} along which a descending heavy body recedes
uniformly, in equal times, from a given point, or approaches it${}^{(a)}$: which so pleased its most
praised Author that he not only several times challenged Friends and Adversaries alike, in the
\emph{Acta}, to attempt it [See the \emph{Yr.}~1689, \emph{Apr.\ p.}~198, and 1690, \emph{May p.}~229,
and \emph{Jul.\ p.}~360,] but himself also toiled strenuously at it, as witnessed by several
properties of the Curve which he had had inserted in the \emph{French Journal}, as I have from the
report of my \emph{Brother}, who nevertheless could not name them to me. My \emph{Brother} himself,
too, was much occupied with the same thing while he was staying at \emph{Paris}, but with all his
inverse Method of Tangents he accomplished nothing else than to reduce the Problem to this
differential equation, $(x\,dx + y\,dy)\sqrt{y} = (x\,dy - y\,dx)\sqrt{a}$, to which, however,

\textbf{(a)}~Leibniz called this the Paracentric Isochrone.

\origpage{602}
one arrives even with no labour${}^{(b)}$. But its full solution, or construction, neither he nor
anyone else was able to give. Recently, while I was preparing the preceding essay for
\emph{Leipzig}, and had already committed to paper part of the discovery together with the figures
themselves, I undertook the scrutiny of this Line, on the occasion of the new Theorems which I gave
there for the radii of the osculating circles in Spirals [whose investigation, however, was itself
made only incidentally; for the chief points of the discovery were only coming to birth under my
pen,] and soon, after a few attempts, I attained the wished-for end${}^{(c)}$, whereupon this most
singular thing befell me, that, the analysis being

\textbf{(b)}~Let $A\theta\omega\alpha$ be the required curve, and let $A\gamma = x$, $\gamma\alpha = y$,
and therefore $A\alpha = \sqrt{x^{2}+y^{2}}$, and its differential
$\alpha\beta = (x\,dx + y\,dy) : \sqrt{x^{2}+y^{2}}$, which is the momentary recession of the body
from the point $A$, while, namely, it describes the little space
$\delta\alpha = \sqrt{dx^{2}+dy^{2}}$, just as, in traversing $A\omega$ the little beginning of the
curve, it recedes from $A$ by the length $A\omega$ itself. Let those recessions $\alpha\beta$,
$A\omega$ be equal, and therefore, by the condition of the curve, the times in which $A\omega$,
$\delta\alpha$ are traversed are equal; consequently these little spaces $A\omega$, or
$\alpha\beta\ [(x\,dx + y\,dy) : \sqrt{x^{2}+y^{2}}]$, and $\alpha\delta\ [\sqrt{dx^{2}+dy^{2}}]$ are
proportional to the velocities of the body at $A$ and $\alpha$. But these velocities, by the law of
descent of heavy bodies, are as the roots of the heights from which the heavy body descended, that
is, as the roots of the heights $iA\ [a]$, from which the heavy body descended at $A$, and
$iA + \alpha\gamma$, $[a+y]$ from which the heavy body descended at $\alpha$; we have therefore this
equation $(a+y)(x\,dx + y\,dy)^{2} : (xx+yy) = a(dx^{2}+dy^{2})$, or
$ax^{2}\,dx^{2} + 2ayx\,dy\,dx + ayy\,dy^{2} + yxx\,dx^{2} + 2yyx\,dy\,dx + y^{3}\,dy^{2} =
axx\,dx^{2} + axx\,dy^{2} + ayy\,dx^{2} + ayy\,dy^{2}$, which, the equal terms being removed from
both sides, the term $2ayx\,dy\,dx$ being transposed, and the square root being extracted, is
reduced to this $(x\,dx + y\,dy)\sqrt{y} = (y\,dx - x\,dy)\sqrt{a}$; which, however, on account of
the mixing together of the indeterminates, can scarcely be constructed.

\textbf{(c)}~Let therefore more convenient indeterminates be chosen, by which, once found, the points
of the curve may nevertheless be determined. Let, for example, $A\alpha = t$, $\varepsilon\zeta = z$,
and it will be $\alpha\gamma = tz : a$, because $A\varepsilon : A\alpha = \varepsilon\zeta : \alpha\gamma$,
and also $\alpha\beta = dt$, and $\beta\delta = t\,dz : \sqrt{aa-zz}$, since
$A\zeta\ [\sqrt{aa-zz}] : A\varepsilon\ [a] = \varepsilon\varkappa\ [dz] :
\varepsilon\lambda\ [a\,dz : \sqrt{aa-zz}]$, and
$A\varepsilon\ [a] : \varepsilon\lambda\ [a\,dz : \sqrt{aa-zz}] = A\delta$ or
$A\alpha\ [t] : \beta\delta\ [t\,dz : \sqrt{aa-zz}]$. Therefore, since in the note above it was shown
that $\alpha\beta : \alpha\delta = \sqrt{iA} : \sqrt{iA + \alpha\gamma}$, or
$\alpha\beta^{2} : \alpha\delta^{2} = iA : iA + \alpha\gamma$, or, by inversion,
$\alpha\beta^{2} : \beta\delta^{2} = iA : \alpha\gamma$, it will be
$dt^{2} : \dfrac{tt\,dz^{2}}{aa-zz} = a : \dfrac{tz}{a}$, which is reduced to
$dt : \sqrt{t} = a\,dz : \sqrt{aaz - z^{3}}$, where now the indeterminates are separated, and,
integrating further, $\displaystyle 2\sqrt{t} = \int(a\,dz : \sqrt{aaz - z^{3}})$, whence it is clear
that the matter has been reduced to the quadrature of a curve whose abscissa is $z$, ordinate
$a^{2}\sqrt{a} : \sqrt{aaz - z^{3}}$.

\origpage{603}
fully worked out, I noticed that the Elastic curve, which had given only some chance occasion to
this attempt, was that very one by whose rectification the other could be constructed${}^{(d)}$. For
although the same might be accomplished by means of the quadrature of some algebraic space, I
nevertheless judge the other mode of construction to be preferred, both because in general curves are
more easily rectified in practice than spaces are squared, and especially because nature herself
[provided she observes a suitable law of tension] seems to have prescribed that one.

\textbf{Construction.} The Elastic curve of the third construction [\emph{Fig.}~3.] having therefore been described, and the
rest being placed as before, let there be applied, in each quadrant of the circle $BL$, $LG$, the
straight line $\varepsilon\zeta$ equal to $Ag$ itself, that is, to the third proportional to $AB$
and $AE$ [here it appears drawn only in the right quadrant, so that confusion of the lines may be
avoided] and from the drawn $A\varepsilon$ [produced, if there be need] let there be cut off
$A\alpha$ equal to the third proportional to the straight line $AB$ and the portion of the curve
$AQ$ in the left, or $\varphi RQ$ in the right quadrant; and the point $\alpha$ will be on a certain
curve $A\chi\omega\eta$, so contrived that a heavy body carried along it [if it has first fallen from
the height $iA$] approaches the point $A$ uniformly in equal times, or recedes from it. The
demonstration of this construction [so that I may relieve the most Celebrated Author of the Problem,
most busied with other matters, of the labour of examining it] I here append.

\rmfigure{figures/fig-3.png}{Fig.~3.}{Fig. 3 (from plate Tab. XXIV, facing p. 618 of the Opera,
Tomus I): the elastic curve, the isochronal curve $A\chi\omega\eta$ and the auxiliary circle,
quadrants and construction lines used in the demonstration.}

\textbf{(d)}~But if, for the sake of a further reduction, one puts $z = uu : a$, it will be
$a\,dz : \sqrt{aaz - z^{3}} = 2au\,du : a\sqrt{au^{2} - u^{6} : a^{3}} =
2aa\,du : \sqrt{a^{4}-u^{4}}\sqrt{a}$. Therefore
$\displaystyle 2\sqrt{t} = \frac{2}{\sqrt{a}}\int(aa\,du : \sqrt{a^{4}-u^{4}})$, or
$\displaystyle \sqrt{at} = \int(aa\,du : \sqrt{a^{4}-u^{4}})$. Now [preceding No., Note~q.]
$\displaystyle\int(aa\,du : \sqrt{a^{4}-u^{4}}) ={}$ the arc $AQ$ or $\varphi RQ$ of the elastic
curve, provided of course $AE = u$ is made. Therefore $\sqrt{at}\ [\sqrt{a\cdot A\alpha}] = AQ$ or
$\varphi RQ$, and $A\alpha = AQ^{2} : a$ or $\varphi RQ^{2} : a$, the third proportional to $AB$,
$AQ$ or $\varphi RQ$, when $\varepsilon\zeta\ [z = uu : a]$ is taken as the third proportional to
$AB$, $AE$. Whence the Author's construction follows. See the following No., Note (a).

\origpage{604}
\textbf{Demonstration.} Let a circle $\beta\delta$ be understood as described with centre $A$, cutting off from the applied
line $A\alpha$ the Element $\alpha\beta$, and from the curve the Element $\alpha\delta$, at whose
vertex let the isochronous little portion $A\omega$ be conceived. Now the ratio of $\beta\delta$ to
$\alpha\beta$ is compounded of the following five ratios:

\[
\begin{aligned}
\beta\delta : \text{Elem. periph. } G\varepsilon\ [\varepsilon\lambda] &= A\alpha : A\varepsilon \\
\text{Elem. per. } G\varepsilon\ [\varepsilon\lambda] : \text{Elem. } \varepsilon\zeta\ [\varepsilon\varkappa]
   &= A\varepsilon : A\zeta, \qquad \text{ex nat. Circuli.} \\
\text{Elem. } \varepsilon\zeta\ [\varepsilon\varkappa] : \text{Elem. } AE
   &= 2AE : AB, \qquad \text{per constr. \& princ. Calc. inf. parvor.} \\
\text{Elem. } AE : \text{Elem. } AQ\ [\varphi RQ]
   &= A\zeta : AB, \qquad \text{ex nat. curvae } AQ\ \text{ut collig. ex Cor. 2. Constr. III. praeced.} \\
\text{Elem. } AQ\ [\varphi RQ] : \alpha\beta
   &= AB : 2AQ\ [2\varphi RQ], \qquad \text{per constr. \& princ. Calc. inf. parvor.}
\end{aligned}
\]

whence, the division being made by the common antecedents and consequents of the ratios,
$\beta\delta : \alpha\beta = A\alpha : AB + AE : AQ\ [\varphi RQ,]$ and
$\beta\delta^{2} : \alpha\beta^{2} = A\alpha^{2} : AB^{2} + AE^{2} : AQ^{2}$ or $\varphi RQ^{2}$
[$\varepsilon\zeta : A\alpha$, by constr.]
$= A\alpha^{2} \times \varepsilon\zeta : AB^{2} \times A\alpha =
A\alpha \times \varepsilon\zeta : AB^{2} ={}$ [on account of the similarity of the Triangles
$A\alpha\gamma$, $A\varepsilon\zeta$] $A\varepsilon \times \alpha\gamma : AB^{2} = \alpha\gamma : AB\ [Ai]$,
and, by composition, $\alpha\delta^{2} : \alpha\beta^{2} = \alpha\gamma + Ai : Ai$; and accordingly
$\alpha\delta : \alpha\beta = \sqrt{\alpha\gamma + Ai} : \sqrt{Ai} ={}$ the velocity acquired at
$\alpha$ : the velocity acquired at $A ={}$ [from the nature of the descent of heavy bodies]
$= \alpha\delta : A\omega$ [by hypothesis, since they are supposed to be traversed in equal little
times]. Therefore $\alpha\beta = A\omega$. Therefore a heavy body carried along this Curve recedes
uniformly, in equal times, from the point $A$. Q.~E.~D.

The things chiefly to be observed here are the following:

\textbf{1.}~If the heavy body be supposed to fall from higher or lower than from $i$, say from
$\tau$, there is no need of a new Elastic curve: this one alone suffices for describing all the
isochronous Curves. For they are all similar to one another, and are determined merely by making, as
$Ai$ to $A\tau$, so $A\alpha$ to another $A\alpha$ to be cut off from the same $A\varepsilon$. And
the curve of our Figure is itself constructed for a heavy body falling as though from $\tau$, not
from $i$.

\textbf{2.}~But there are also infinitely many others which a heavy body, having slipped down from
the same height

\origpage{605}
$\tau A$, can traverse so as to approach or recede uniformly from some point; but that point
continually different from $A$; nor, except by a single curve from the same height, with respect to
the same point, can uniform approach or recession be brought about.${}^{(e)}$

\textbf{3.}~Our Isochronous Curve takes its beginning at the very point $A$, with respect to which
the motion is uniform, and is terminated at the point $\eta$ of the same height as it; at each
extremity it touches the horizontal straight line $A\eta\,{}^{(f)}$. Whence it is established to have
a contrary flexure. But if on the other side there be set up a Curve $A\lambda\theta\psi$ similar and
equal to this, and the heavy body be released from $\eta$ with the velocity which it can acquire by
falling from a height equal to $\tau A$ itself, then, first by descending to $\omega$, then by
re-ascending, it will approach $A$ uniformly, and thence, going on continually through the other
$A\lambda\theta\psi$, it will be drawn away from the same point $A$ by the same law.

\textbf{4.}~The tangent of the curve at intermediate places, such as $\alpha$, is found if, having
drawn $A\mu$ perpendicular to $A\alpha$, one makes, as $\sqrt{A\tau}$ to $\sqrt{\alpha\gamma}$, so
$A\alpha$ to $A\mu$; for $\alpha\mu$ being drawn will touch the curve: which is easily inferred even
from the use which the curve ought to serve${}^{(g)}$.

\textbf{5.}~The straight line $A\eta$ is quadruple the straight line $A\theta$.${}^{(h)}$

\textbf{6.}~If on the Elastic curve a point $Q$ be taken such that the third proportional to the
subtangent $Pp$ and the straight line $AB$ equals the curve $AQ$, there will be found

\textbf{(e)}~Indeed there are infinitely many, as Leibniz observed in No.~LXIV, and Huygens in
No.~LXV, and the Author acknowledged in No.~LXVI. Namely, in the integration of the equation
$dt : \sqrt{t} = a\,dz : \sqrt{aaz - z^{3}}$, [Note (c)] the addition of the constant was omitted.
But the complete integral $\displaystyle 2\sqrt{t} + \sqrt{b} = \int(a\,dz : \sqrt{aaz - z^{3}})$
pertains to infinitely many different curves, according as one quantity $b$ or another is assumed.
See No.~LXVI.

\textbf{(f)}~See the following Note.

\textbf{(g)}~For $A\alpha : A\mu = \alpha\beta : \beta\delta ={}$ [See Note c] $\sqrt{iA}$ or here
$\sqrt{\tau A} : \sqrt{\alpha\gamma}$. Therefore if $\alpha\gamma = o$, which takes place where
$A\alpha$ falls on $AG$, then $A\mu$ too will be infinitely less than $A\alpha$ itself, and so the
curve touches the straight line $A\alpha$, that is $AG$. Of this kind are the points $A$ and $\eta$,
which Our [Author] called the beginning and end of the curve, although it is in truth unterminated,
and has infinitely many flexures and windings. See again No.~LXVI.

\textbf{(h)}~For $A\eta = AR\varphi^{2} : AB$ [by constr.], but $A\theta$ $= AQR^{2} : AB$. Therefore
since $AR\varphi = 2AQR$ it will be $A\eta = 4A\theta$.

\origpage{606}
by its means the point $\chi$ of the portion $A\chi\theta$ which of all is the most distant from the
perpendicular $A\theta\,{}^{(i)}$.

\textbf{7.}~But if the point $Q$ be placed such that the tangent $Qp$ equals the curve $\varphi RQ$,
by its means there will be had, on the other curve, the point $\omega$, the lowest of all, and the
most remote from the horizontal $A\eta\,{}^{(k)}$.

\textbf{8.}~Finally, this too is to be noted, that from the construction of the proposed Problem it
is now easy to gather what suppositions are to be made in order to reduce the equation
$(x\,dy + y\,dx)\sqrt{y} = (y\,dx - x\,dy)\sqrt{a}$. For if in place of $x$ and $y\,{}^{(l)}$ there be
assumed two other indeterminates $t$ and $z\,{}^{(m)}$, by putting $ay = tz$, and
$ax = t\sqrt{aa - zz}$; there will come forth an equation in which the indeterminates $t$ and $z$,
with their differentials, can be separated from one another; which suffices for then expediting the
construction, at least by quadratures. For two things in this calculus seem still chiefly to be
wanting, which, if they could always be done, everything would be found out; one, that differentials
of the second or other kinds be reduced to differentials of the first;

\textbf{(i)}~The tangent at the point $\chi$, which is the most remote from the Axis $AL$, is
parallel to the straight line $\alpha\gamma$, and gives the Triangle $A\alpha\mu$ similar to the
Triangle $A\alpha\gamma$. Therefore $\alpha\gamma^{2} : A\gamma^{2} = A\alpha^{2} : A\mu^{2} ={}$
[Note~g] $Ai : \alpha\gamma$, or, by composition, $\alpha\gamma^{2}\ [ttzz : aa] : A\alpha^{2}\ [tt] = \uncertain{Ai\ [a] :}
Ai + \alpha\gamma\ [a + tz : a]$; whence is elicited $t\ [AQ^{2} : a] = (a^{4} - aazz) : z^{3}$ [since
$z = uu : a$] $= (a^{4} - u^{4})\,a^{3} : u^{6} = a^{3} : Pp^{2}$. For [preceding No., Note (c)]
$Pp = u^{3} : \sqrt{a^{4}-u^{4}}$ and therefore $Pp^{2} = u^{6} : (a^{4}-u^{4})$. Therefore since at
the point $\chi$, $AQ^{2} : a = a^{3} : Pp^{2}$, or $AQ^{2} = a^{4} : Pp^{2}$, it is also
$AQ = aa : Pp$, that is, $AQ$ the third proportional to the tangent $Pp$ and the straight line $AB$.

\textbf{(k)}~But the tangent at the point $\omega$, which is most distant from the Axis $AG$, is
parallel to it and makes the Triangle $A\alpha\mu$ similar to the Triangle $A\alpha\gamma$. Therefore
$A\gamma^{2} : \alpha\gamma^{2} = A\alpha^{2} : A\mu^{2} = Ai : \alpha\gamma$, or, by composition,
$A\alpha^{2}\ [tt] : \alpha\gamma^{2}\ [ttzz : aa] = Ai + \alpha\gamma\ [a + tz : a] :
\alpha\gamma\ [tz : a]$; whence is deduced $t = [\varphi RQ^{2} : a] = aaz : (aa - zz) =
a^{3}uu : (a^{4} - u^{4})$ [namely, by writing $uu : a$ for $z$] $= Qp^{2} : a$. For
[preceding No., Note (c)] $Qp = aau : \sqrt{a^{4}-u^{4}}$. Therefore at the point $\omega$,
$\varphi RQ^{2} : a = Qp^{2} : a$, or the arc $\varphi RQ ={}$ the tangent $Qp$.

\textbf{(l)}~$A\gamma$ and $\gamma\alpha$.

\textbf{(m)}~$A\alpha$, and $\varepsilon\zeta$.

\origpage{607}
the other, that in differential equations of the first kind the indeterminates, if they be mixed
together, be mutually separated from one another, so that each with its own differential may
constitute a particular part of the equation; that is, that the inverse method of tangents be found.
For each I have given certain rules [and I could give more without end], similar to those which my
\emph{Brother} deposited at \emph{Paris} with the \emph{Marquis} de l'Hospital, in which he had at
first supposed a method to reside. But I at once perceived that they contained nothing but certain
particular artifices, which I would not dare to call a method; inasmuch as I judge that no such thing
can be given, any more than a universal method can be given for constructing algebraic equations of
any degrees whatsoever indiscriminately. For that one [of which mention was once made in the
\emph{Acta}], by which I posit indefinitely $o = a + bx + cy + exy +{}$~etc.\ and then seek by its
means the tangent of the line, seems scarcely to succeed otherwise than in speculation. That this
may lend credence to what has been said, let it be this: that the present Problem, although as it
appears not very difficult, could be solved by no one by the aid of any rule or method hitherto
discovered.

\end{document}
